% Answers to Chapter 4, translated from sv/back/facit/04.tex.

\facitavsnitt{c4}{Chapter 4}{Systems of equations}
\begin{facit}

\svar{1.1} \sv{a} $9.4\,\text{V}$ \sv{b} $1.88\,\text{W}$

\svar{1.2} \sv{a} $R = \frac{\tau}{C} = 2\,000\,\Omega$
\sv{b} $2\,\text{k}\Omega$

\svar{1.3} $5U_{\text{in}} \leq 15$, so $U_{\text{in}} \leq 3\,\text{V}$.

\svar{2.1} $x = \frac{20}{7} \approx 2.86$, $y = \frac{17}{7} \approx 2.43$

\svar{2.2} \sv{a} $I_1 + I_2 = 5$ and $I_1 = 2I_2$
\sv{b} $I_1 = \frac{10}{3} \approx 3.33\,\text{A}$, $I_2 = \frac{5}{3} \approx 1.67\,\text{A}$

\svar{2.3} \sv{a} $U_1 + U_2 = 12$ and $U_2 = 3U_1$
\sv{b} $U_1 = 3\,\text{V}$, $U_2 = 9\,\text{V}$
\sv{c} $R_1 = 1.5\,\Omega$, $R_2 = 4.5\,\Omega$

\svar{3.1} \sv{a} $(2, 5)$ \sv{b} $(9, 3)$ \sv{c} $(7, 3)$

\svar{3.2} \sv{a} $(3, 1)$ \sv{b} $(4, 2)$ \sv{c} $(3, 2)$ \sv{d} $(2, 1)$

\svar{3.3} \sv{a} Infinitely many: equation (2) is equation (1) times $2$, so the lines coincide.
\sv{b} None: the same left-hand side but different right-hand sides, so the lines are parallel.
\sv{c} Exactly one, $(2, 2)$: the lines have different slopes and intersect.
\sv{d} Infinitely many: equation (2) times $-2$ is equation (1), so the lines coincide.

\svar{3.4} \sv{a} For $x = 0, 1, 2, 3, 4$, $y = 2x - 1$ gives the values $-1, 1, 3, 5, 7$ and
$y = -x + 5$ the values $5, 4, 3, 2, 1$.
\sv{b} At $x = 2$; the point of intersection is $(2, 3)$.
\sv{c} $(x, y) = (2, 3)$, the same as the reading.
\sv{d} The rows would never give the same $y$: the difference between them would be constant, as
for two parallel lines.

\svar{3.5} \sv{a} $(4, 2)$ \sv{b} $(6, 3)$

\svar{3.6} \sv{a} $I_1 + I_2 = 12$ and $I_1 - I_2 = 4$
\sv{b} $I_1 = 8\,\text{A}$, $I_2 = 4\,\text{A}$
\sv{c} $8 + 4 = 12$ and $8 - 4 = 4$ both hold.

\svar{3.7} \sv{a} $I_1 + I_2 = 3$ and $I_1 = 2I_2$
\sv{b} $I_1 = 2\,\text{A}$, $I_2 = 1\,\text{A}$
\sv{c} $R_1 = 12\,\Omega$, $R_2 = 24\,\Omega$
\sv{d} $R_1 \parallell R_2 = \frac{12 \times 24}{12 + 24} = 8\,\Omega = \frac{24}{3}\,\Omega$

\svar{3.8} \sv{a} $U_1 + U_2 = 15$ and $U_1 - U_2 = 3$
\sv{b} $U_1 = 9\,\text{V}$, $U_2 = 6\,\text{V}$
\sv{c} $R_1 = 30\,\Omega$, $R_2 = 20\,\Omega$
\sv{d} $R_1 + R_2 = 50\,\Omega = \frac{15}{0.3}\,\Omega$

\svar{3.9} \sv{a} $U = 12 - 2I$ and $U = 4I$
\sv{b} $I = 2\,\text{A}$, $U = 8\,\text{V}$
\sv{c} $12 - 2 \times 2 = 8$ and $4 \times 2 = 8$ both hold.
\sv{d} The operating point.

\svar{3.10} \sv{a} $\frac{6}{9} = \frac{R_1}{900}$, which gives $R_1 = 600\,\Omega$.
\sv{b} $R_1 = 600\,\Omega$, $R_2 = 300\,\Omega$
\sv{c} $U_1 = 9 \times \frac{600}{900} = 6\,\text{V}$, which is correct.

\svar{3.11} \sv{a} $I_1 = 3\,\text{A}$, $I_2 = 2\,\text{A}$, $I_3 = 4\,\text{A}$
\sv{b} $3 + 2 + 4 = 9$, $3 - 2 = 1$ and $I_3 = 4$ all hold.
\sv{c} It gives one unknown directly; substituted into the others, it leaves a system with two
unknowns.

\svar{3.12} \sv{a} $5R_A + 3R_B = 4\,390$ and $2R_A + 4R_B = 3\,660$
\sv{b} $R_A = 470\,\Omega$, $R_B = 680\,\Omega$
\sv{c} $2\,350 + 2\,040 = 4\,390$ and $940 + 2\,720 = 3\,660$ both hold.

\svar{3.13} \sv{a} Yes. \sv{b} Yes. \sv{c} No: $5 - 2 \times 1 = 3$ holds, but
$5 + 1 = 6 \neq 7$.
\sv{d} The solution must lie on both lines at once. A pair of numbers that satisfies only one of
the equations lies on only one of the lines, as in c).

\svar{3.14} \sv{a} $(2, 3)$ \sv{b} $(4, -1)$ \sv{c} $(2, -1)$
\sv{d} Substitution in a) and b), where one unknown is easy to solve for; elimination in c), where
the $y$-terms easily become opposite. Both methods give the same answer.

\svar[Test yourself]{T} \sv{1} $(x, y) = (3, 2)$
\sv{2} $I_1 = 3\,\text{A}$, $I_2 = 1\,\text{A}$

\end{facit}
